%!TEX root = ../../../exa-ma-d7.1.tex
\section{PC5 demonstrator context}
\label{sec:pc5:demonstrator-context}

PC5 demonstrators integrate components and methods developed across the
scientific work packages. An application benchmark may exercise one of
these components and provide evidence for its numerical behaviour or
performance within the demonstrator. The scope of each reported result
remains that of the component and configuration actually measured.

\subsection{Ktirio Urban Building}
\label{sec:pc5:ktirio-urban-building}

Ktirio Urban Building addresses urban building energy simulation and
decision support, in joint work with the HiDALGO2 Centre of Excellence.
The demonstrator connects terrain and building geometry, urban microclimate
and building energy models, and ensemble uncertainty quantification.
The following Exa-MA contributions support these roles.

\paragraph{Geometry generation and its benchmark application.}
Ktirio-Geom is the component that generates terrain and building geometry
from heterogeneous geographic inputs. Its developments include roof
construction, tiled terrain generation, surface remeshing and volume
meshing. The application \texttt{app-feelpp-ktirio-geom}
(\Cref{sec:app:specs:app-feelpp-ktirio-geom}) benchmarks this component;
\texttt{BM-FEEL-KTIRIO-01} reports the terrain-generation campaign.
These contributions connect WP1/T1.1 meshing methods with WP7 implementation
and benchmarking. The existing terrain timings cover one component of the
urban demonstrator.

\paragraph{Distributed geometric queries.}
The Feel++ two-level BVH and asynchronous MPI queries provide distributed
geometric-search infrastructure for the demonstrator. This WP1/T1.4
contribution complements Ktirio-Geom's geometry generation; its software
implementation is documented in the Feel++ fiche.

\paragraph{Urban microclimate and building energy coupling.}
The preCICE integration targets coupling urban microclimate with building
energy simulation while allowing the solvers to retain their own
interface meshes and discretizations. WP1/T1.5 describes interface
field transfer, with a connection to the multiphysics context of T1.7.
WP7 documents the Feel++ adapter and OpenFOAM participant deployment.
The same infrastructure also supports broader Feel++--OpenFOAM coupling
objectives.

\paragraph{Ensemble UQ and decision support.}
Large ensembles of urban building energy simulations create memory and
I/O challenges through their building-level outputs and time series.
The Ensemble/UQ framework resides in Ktirio Urban Building and builds on
Feel++ streaming statistical processing. The WP6/T6.4 contribution, with
relevance to T6.2, concerns that statistical-processing layer: mergeable
summaries and quantiles provided by feelstats. Ensemble execution,
dedicated I/O ranks and application-level VaR/CVaR risk products belong
to the Ktirio framework. The Feel++ fiche describes the statistical
components and their application context. Geometry benchmark results and
energy/UQ assessments retain their respective scopes within this common
demonstrator.
